\documentclass[11pt]{article}
\usepackage[margin=1in]{geometry}
\usepackage[T1]{fontenc}
\usepackage{lmodern}
\usepackage{amsmath,amssymb}
\usepackage{booktabs}
\usepackage{graphicx}
\usepackage{xcolor}
\usepackage{hyperref}
\usepackage{url}
\usepackage{caption}
\hypersetup{colorlinks=true,linkcolor=blue!50!black,citecolor=blue!50!black,urlcolor=blue!50!black}

\newcommand{\figplaceholder}[2]{\fbox{\parbox[c][#2][c]{0.95\linewidth}{\centering\small\color{gray} #1}}}

\title{osvcut: Open, Scriptable Processing of Consumer Dual-Fisheye 360\textdegree{} Video\\with Flow-Aligned Minimum-Disagreement Seams}
\author{Author(s) to be confirmed}
\date{Draft --- \today}

\begin{document}
\maketitle

\begin{abstract}
Consumer 360\textdegree{} cameras record two fisheye streams that must be stitched, levelled and reframed before footage can be edited, and the vendor tools that do this are graphical and closed.
We present \emph{osvcut}, an open pipeline that turns DJI Osmo~360 \texttt{.OSV} recordings into ordinary, scriptable video: a schema-free decoder for the per-frame metadata, horizon levelling from the embedded IMU with an empirically recovered axis map, a GPU renderer driven by a JSON edit specification, and a browser studio for people and AI agents.
On an 8.1-hour personal corpus (239 clips) we analyse why calibrated dual-fisheye stitches still show seams: near-field parallax, inter-lens exposure differences and a body-occluded sector where the two lenses barely overlap; we rule out vignetting, rolling-shutter mismatch (with a per-pixel model), systematic calibration error and a gravity-aligned ground-plane parallax model.
We then propose a seam that aligns both lenses half-way with dense optical flow and places a narrow morph transition on a minimum-disagreement closed path through the overlap.
To compare seam methods fairly we introduce a method-agnostic \emph{contribution disagreement} metric and evaluate on 100 frames from 20 clips: the proposed seam reduces disagreement by 38\% relative to the feathered stitch and by 11\% relative to a graph-style minimum-difference cut (79/100 paired wins, $p<10^{-12}$), and by 9\% on the parallax-heavy subset, without adding temporal flicker relative to a no-seam reference (9 clips).
Code is released under Apache-2.0.
\end{abstract}

% ---------------------------------------------------------------------------
\section{Introduction}

Pocket 360\textdegree{} cameras have made spherical video a casual medium, but their footage is not editable as recorded.
A DJI Osmo~360 clip stores two 10-bit HEVC fisheye streams, one per lens, plus calibration and inertial metadata in a proprietary protobuf track.
To obtain a normal video one must stitch the fisheyes into a sphere, compensate the camera's attitude and then aim a virtual camera (\emph{reframing}).
Vendor software does this interactively; existing open alternatives integrate into commercial editors~\cite{openosv}.
Neither exposes the footage to scripts, which is what automated or AI-assisted editing requires: an agent needs to read what is in a clip, decide where to look and when to cut, and render the result reproducibly.

This paper describes an open pipeline built for that purpose and the seam research it required.
Our contributions are:
\begin{enumerate}
\item \textbf{Open access to the format.} A schema-free decoder for the \texttt{djmd} metadata (lens calibration, per-frame attitude) that reads a whole clip by stream copy in seconds, and an empirically recovered map from IMU to lens axes that enables horizon levelling from metadata alone (\S\ref{sec:format}).
\item \textbf{An analysis of seam failures} in calibrated dual-fisheye stitching on a real corpus, including three negative results (vignetting, rolling shutter, ground-plane parallax) (\S\ref{sec:why}).
\item \textbf{A seam method} combining lens-validity masks, a per-azimuth exposure field and flow-aligned minimum-disagreement transitions (\S\ref{sec:seam}).
\item \textbf{An evaluation protocol} with a method-agnostic disagreement metric, a parallax-heavy subset and a no-seam reference for temporal stability, plus a blind pairwise study tool (\S\ref{sec:eval}).
\item \textbf{A scriptable editing system}: a JSON render specification, AI-readable contact sheets, IMU motion and speech timelines, subject following and multi-clip sequences, exposed through an HTTP API and a browser studio (\S\ref{sec:system}).
\end{enumerate}

% ---------------------------------------------------------------------------
\section{Related work}

\paragraph{Dual-fisheye stitching.} Ho and Budagavi~\cite{ho2017dualfisheye} stitch dual-fisheye images by unwarping them, compensating light fall-off and refining the alignment of the overlap region. Lens models for wide fields of view follow Kannala and Brandt~\cite{kannala2006generic}, which DJI's calibration also uses.
\paragraph{Seams and blending.} Optimal seams through overlaps are classically found with dynamic programming or graph cuts~\cite{kwatra2003graphcut,agarwala2004photomontage}, and blended with multi-band splines~\cite{burt1983multiresolution}. Parallax-tolerant methods warp images locally~\cite{zaragoza2013apap,zhang2014parallax}.
\paragraph{Flow-based panoramic video.} Camera-array systems interpolate between views with optical flow to hide parallax~\cite{anderson2016jump,perazzi2015panoramic}. Our setting differs: two back-to-back lenses with a very small baseline, a fixed overlap ring a few degrees wide that collapses to about one degree on the body-occluded side, and per-frame IMU data.
\paragraph{Gyro-based stabilisation.} Gyroscope data has long been used for video stabilisation and rolling-shutter correction~\cite{karpenko2011gyro}; we use the camera's fused attitude for levelling and test whether rolling shutter affects the seam.

% ---------------------------------------------------------------------------
\section{The OSV format and its metadata}\label{sec:format}

An \texttt{.OSV} file is an ISO base media file. ffmpeg demuxes two HEVC Main10 video tracks (\emph{slave}, optical axis $-Y$; \emph{master}, $+Y$), AAC audio, two \texttt{djmd} metadata tracks and two debug tracks.
In our corpus, 226 clips use 4K mode (two $1920^2$ streams at 100\,fps), 12 use 8K mode (two $3840^2$ streams at 29.97\,fps, no audio).

\paragraph{Calibration.} Sample~0 of the primary \texttt{djmd} track carries 24 calibration slots; each holds $f_x,f_y,c_x,c_y$, five radial coefficients $k_1..k_5$ of the Kannala--Brandt model in a $3840^2$ calibration space, an extrinsic quaternion and a 14-point polygon outlining the camera body seen at the bottom of each image. A ray at angle $\theta$ from the optical axis and azimuth $\phi$ projects to
\begin{equation}
r(\theta)=\theta\,(1+k_1\theta^2+k_2\theta^4+k_3\theta^6+k_4\theta^8+k_5\theta^{10}),\quad
(u,v)=\tfrac{W}{2}+s\big(f_x r\cos\phi + c_x-1920,\; f_y r\sin\phi + c_y-1920\big),
\end{equation}
with stream scale $s=f_{\mathrm{digital}}/f_x$ (0.5065 in 4K mode).

\paragraph{Fast metadata access.} Dumping packets with ffprobe is slow for long clips. Because every sample is a \texttt{ProductMeta} message whose later samples contain only field~3 (\texttt{FrameMeta}), the stream-copied track parsed as \emph{one} protobuf message is a flat list in which each frame contributes exactly one field~3. Reading attitudes for a 19-minute clip (115\,443 frames) takes 22\,s instead of minutes.

\paragraph{IMU-to-lens axes.} Each frame carries a fused attitude quaternion $q$ and an accelerometer vector. $R(q)^\top\hat z=(0.639,0.321,0.699)$ matched the normalised negative acceleration $(0.639,0.367,0.676)$, so $R(q)$ maps IMU to world. Rendering a levelled panorama for all 24 signed axis permutations of the up vector left exactly one with vertical structures: $\mathbf{b}=(-\iota_y,-\iota_x,\iota_z)$ for IMU coordinates $\iota$. Levelling then rotates each frame so that world up is the panorama's up while the heading follows a 1-s moving average of the camera's forward direction.

% ---------------------------------------------------------------------------
\section{Why calibrated stitches still show seams}\label{sec:why}

We unwrap both lenses along the seam band (latitude $\beta\in[-8^\circ,8^\circ]$ around the plane between the lenses, $0.1^\circ$/px) and compare them.

\begin{itemize}
\item \textbf{Calibration is correct at infinity.} Distant ridges and clouds coincide.
\item \textbf{Parallax.} Misalignment has a median of $0.77$--$0.85^\circ$ and reaches $2.7^\circ$ on objects within one metre (the hand on the selfie stick, people). The calibration encodes rotation only.
\item \textbf{Exposure.} The slave/master luminance ratio varies from $0.53$ to $1.40$ around the ring (strong backlight).
\item \textbf{Occlusion.} The body polygon marks where the camera itself is visible. Both lenses' occluded sectors face the handle, so the usable overlap shrinks from about $10^\circ$ to about $1^\circ$ there, and the occlusion edge is a $1$--$2^\circ$ dark gradient that a naive blend mixes in.
\end{itemize}

\paragraph{Ruled out.} (i) \emph{Vignetting}: the median radial luminance profile is flat ($\pm4\%$) to $96^\circ$. (ii) \emph{Rolling shutter}: over 40 frames with angular speeds of $5$--$533^\circ$/s, seam flow does not grow with speed (Pearson $r=0.09$). Because a rolling-shutter displacement depends on azimuth (a seam point may be read near the top of one sensor and the bottom of the other) and could cancel in such a summary, we also fit a per-pixel model: with readout time $\tau$ and body angular velocity $\omega$ from the IMU, the second lens sees a seam direction $u$ later by $\Delta t=\tau(y_1-y_0)/H$, predicting a displacement $-(\omega\times u)\,\Delta t$. Over $\sim$1.5\,M confident pixels from 30 frames per clip, the fit explains nothing ($R^2\le0.000$, identical to a control with $\omega$ from shuffled frames). (iii) \emph{Calibration error}: per $10^\circ$ azimuth bin, the confident flow has no bias that persists across 30 frames (RMS of the per-bin mean $0.04^\circ$ against a per-frame spread of $0.12^\circ$; no bin with $|t|>4$ on one clip, one on the other), so large local misalignments are frame-specific. (iv) \emph{A gravity-aligned ground plane}: fitting $\Delta\beta = k\sin(-\mathrm{el}) + c$ to confident seam flow below the horizon gives $R^2\le0.07$, with $k$ changing sign between clips. (v) A large share of the remaining lens-to-lens difference on texture is not geometric: even after raw-flow alignment the mean absolute difference on textured pixels falls only from $38.4$ to $34.1$ grey levels (from $36.0$ to $32.0$ with anti-aliased sampling), reflecting edge-of-field blur, chromatic aberration and noise.

The last point motivates our design: the transition should run where the two lenses \emph{agree}, not through the geometric centre of the overlap.

% ---------------------------------------------------------------------------
\section{Seam method}\label{sec:seam}

All seam computations live in the camera body frame, so one seam serves every output projection.

\paragraph{Validity masks.} For each lens we rasterise the image circle at $\theta=96.5^\circ$, intersect it with the coded frame and with the inside of the body-occlusion arc, erode by $1^\circ$ and feather. Rays outside the usable area are clamped to its edge, so a sector seen by neither lens is filled from the nearer lens's edge instead of dark pixels.

\paragraph{Exposure field.} On the overlap we compute, in linear light, the slave/master ratio in $5^\circ$ azimuth bins, fill empty bins circularly, smooth it, and filter it over time. Its circular mean is applied to the whole hemisphere and the azimuthal residual is faded out within $45^\circ$ of the seam.

\paragraph{Flow alignment.} On a seam strip ($\pm4.5^\circ$, $0.2^\circ$/px) we estimate dense flow $F$ with $A(p)\approx B(p+F)$ using Farneb\"ack's method~\cite{farneback2003}, keep it where forward--backward consistency is high, fill elsewhere by normalised convolution and clamp it to $4^\circ$. Dense inverse search~\cite{kroeger2016dis} produced a systematic $+2$--$3.4^\circ$ bias on these 23--46-row strips and was rejected.

\paragraph{Minimum-disagreement path.} We align both strips half-way, $A_{1/2}(p)=A(p-\tfrac12F)$ and $B_{1/2}(p)=B(p+\tfrac12F)$, and find a closed path $\beta^*(\alpha)$ minimising
\begin{equation}
C(\alpha,\beta)=|A_{1/2}-B_{1/2}| + \tfrac14|\nabla A_{1/2}-\nabla B_{1/2}| + 3\,[\text{not both valid}] + 2\,[\text{room}<0.6^\circ] + 0.02\,|\beta-\bar\beta(\alpha)| + \lambda_t|\beta-\beta^*_{t-1}(\alpha)|
\end{equation}
by dynamic programming with slope $\le1$ row per column, starting and ending at the cheapest node so that the path closes around the ring.

\paragraph{Morph transition.} For an output ray with seam coordinates $(\alpha,\beta)$, $t=\mathrm{smoothstep}\big((\beta-\beta^*(\alpha))/2w+\tfrac12\big)$ with half-width $w=\min(1.5^\circ,\text{room in the overlap})$. The slave lens is sampled at $(\alpha,\beta)-tF$, the master at $(\alpha,\beta)+(1-t)F$, and the two are blended with weights $1-t$ and $t$ multiplied by the validity masks. The flow is attenuated where the transition must be narrow.

% ---------------------------------------------------------------------------
\section{Evaluation}\label{sec:eval}

\paragraph{Data.} 20 clips drawn with a fixed seed from the 239 clips longer than 3\,s; 5 evenly spaced frames each ($n=100$). The corpus contains people's faces and is not redistributable; code and metrics are released.

\paragraph{Metrics.} All metrics are computed on the seam strip ($\pm5^\circ$, $0.1^\circ$/px) where both lenses are valid, against single-lens renders of the same directions.
\begin{itemize}
\item \textbf{Contribution disagreement} (CD): $\sum 4w_Aw_B\,|Y_A-Y_B| \,/\, \sum 4w_Aw_B$ over the two (warped, exposure-compensated) contributions, i.e.\ how much the sources disagree wherever they are mixed. It applies equally to crossfades, cut lines and morph ramps.
\item \textbf{Edge deviation}: $|\log(E_\mathrm{out}/E_\mathrm{ref})|$ of across-seam gradient energy within $1.5^\circ$; blur lowers it, a hard visible cut raises it.
\item \textbf{Fidelity}: mean of $\max(\mathrm{SSIM}(\mathrm{out},A),\mathrm{SSIM}(\mathrm{out},B))$~\cite{wang2004ssim}. It favours methods that copy source pixels and is reported for completeness.
\item \textbf{Exposure step} in excess of the natural scene gradient.
\item \textbf{Hard subset}: azimuths where the median lens-to-lens flow exceeds $0.4^\circ$ (dilated by $\pm2^\circ$).
\end{itemize}

\begin{table}[t]
\centering\small
\caption{Seam benchmark, $n=100$ frames from 20 clips; means with 95\% bootstrap intervals for CD. Runtime is update+render of one $3840\times1920$ frame on an RTX 4070 Ti.}
\label{tab:bench}
\begin{tabular}{lcccccr}
\toprule
Method & CD $\downarrow$ & CD hard $\downarrow$ & Edge hard $\downarrow$ & Fidelity hard $\uparrow$ & Step $\downarrow$ & ms \\
\midrule
Feathered ($12^\circ$ crossfade) & 0.101 [.093,.109] & 0.133 & 0.188 & 0.853 & 0.030 & 23 \\
Min-difference cut (DP)          & 0.071 [.066,.075] & 0.092 & \textbf{0.052} & \textbf{0.967} & 0.023 & 146 \\
Flow, ramp centred in overlap    & 0.085 [.079,.091] & 0.119 & 0.076 & 0.870 & 0.028 & 169 \\
Flow + DP path, $\pm3^\circ$     & 0.071 [.066,.076] & 0.097 & 0.086 & 0.884 & \textbf{0.023} & 180 \\
\textbf{Flow + DP path, $\pm1.5^\circ$ (ours)} & \textbf{0.063} [.059,.067] & \textbf{0.084} & 0.065 & 0.929 & 0.025 & 180 \\
\bottomrule
\end{tabular}
\end{table}

\paragraph{Results.} Table~\ref{tab:bench}. Against the minimum-difference cut, our seam has lower CD on 79/100 frames (Wilcoxon $p=1.6\times10^{-13}$) and lower CD on the hard subset on 82/100 ($p=1.7\times10^{-10}$); the cut keeps slightly crisper edges on the hard subset (edge deviation $+0.014$, $p=0.022$) and, by construction, higher fidelity. Against the feathered stitch our seam is better on every metric except the exposure step (n.s., $p=0.25$). Centring a flow morph in the overlap is clearly worse than placing it on the minimum-disagreement path, and a narrow transition beats a wide one.

\paragraph{Temporal stability.} On one test clip (300 frames, crowd within one metre) we measure the mean absolute frame difference inside the seam band relative to a nearby band on the same side. A single-lens render with no seam gives the reference ratio 1.510; the feathered stitch gives 1.155 --- lower than the reference because ghosting low-passes the seam band, not because it is more stable. The minimum-difference cut gives 1.507 and our seam 1.361: neither adds flicker. Across 8 further clips (random 4-s windows, 25\,fps, rendered in memory), the excess over each clip's single-lens reference is $-0.032$ [$-0.272$, $+0.072$] for the cut and $-0.039$ [$-0.200$, $+0.071$] for our seam (3/8 clips above the reference for both), versus $-0.206$ for the feathered stitch.

\paragraph{Perceptual study.} The studio includes a blind pairwise test: a random clip, time and parallax-heavy seam location, rendered as an upright $55^\circ$ close-up by two hidden methods in random order; raters answer left/same/right and scores are fitted with a Bradley--Terry model~\cite{bradley1952rank}. \textcolor{red}{[Results pending: data collection with human raters.]}

\begin{figure}[t]
\centering
\figplaceholder{Qualitative comparison on near-field seams (feathered / min-difference cut / ours). Requires consent for footage showing people, or crops without faces.}{4cm}
\caption{Qualitative comparison (to be added).}
\end{figure}

% ---------------------------------------------------------------------------
\section{A scriptable editing system}\label{sec:system}

\paragraph{Render specification.} A JSON document names the source, in/out times, output projection and size, stabilisation and seam modes, and camera keyframes (yaw, pitch, field of view, and an eye-offset distortion that spans rectilinear to stereographic ``little planet'' views); keyframes interpolate with smoothstep or linearly. The renderer decodes both streams with ffmpeg, stitches and reframes on the GPU and encodes with NVENC.

\paragraph{AI-readable views.} For an agent to decide where to look, each clip can be summarised as levelled panorama contact sheets with yaw ticks, a speech transcript~\cite{radford2023whisper} and an IMU motion curve (angular speed, high-frequency shake, tilt). In a case study, a language model read a 49-s contact sheet and transcript, identified the location and the subjects' bearings, and wrote a keyframe specification; coarse sampling caused the subject to leave the frame in one segment, which motivated subject following.

\paragraph{Subject following.} A user (or agent) draws a box on the 360\textdegree{} view; we render a rectilinear view around the subject each frame with the same levelling as the export, track it with an online multiple-instance tracker~\cite{babenko2009mil}, re-aim the view when the subject leaves the central region, and emit smoothed linear keyframes. On a 5.5-s walking sequence with partial occlusion the subject stayed boxed throughout, at 0.65$\times$ real time.

\paragraph{Live preview.} Opening a clip builds a dual-fisheye proxy (both lenses side by side, decoded and scaled on the GPU at about $4\times$ real time) and a small rig description: per-lens Kannala--Brandt parameters in proxy pixels, the usable radius per image angle, a global exposure balance and the levelling rotation as quaternions at 15\,Hz. A fragment shader evaluates the same lens model per output pixel, so any clip can be browsed, reframed and keyframed in real time without waiting for a stitched preview; exports use the full flow seam.

\paragraph{Studio.} A local web application indexes a folder of clips (239 clips scanned in under 20\,s), shows stitched and levelled panorama thumbnails, generates previews through a GPU job queue, plays them in a WebGL viewer sharing the renderer's projection maths, overlays motion, speech, keyframes and in/out points on the timeline, and assembles multi-clip sequences that are rendered at a common size and frame rate and concatenated losslessly.

% ---------------------------------------------------------------------------
\section{Limitations}

Metrics are proxies: there is no ground truth for a view from the centre of the rig. The perceptual study is pending. The body-occluded sector leaves about one degree of overlap; near objects there remain stretched or cut. Evaluation uses one camera model and one owner's footage. Tracking uses a weight-free tracker and has not been benchmarked against annotations. Runtime ($\approx$0.18\,s per 4K equirect frame) is below real time.

\section{Conclusion}

Consumer dual-fisheye footage can be made as scriptable as ordinary video with open tools. On real footage, seams are dominated by near-field parallax, exposure mismatch and the body-occluded sector; aligning the lenses with optical flow and transitioning along a minimum-disagreement path reduces lens-to-lens disagreement significantly without adding flicker.

\bibliographystyle{plain}
\bibliography{refs}

\end{document}
